%=============================================================
%  PLANTILLA: ARTÍCULO DE REVISIÓN EN COMPUTACIÓN
%  Basada en: Kitchenham & Brereton (2007)
%  Escuela Profesional de Ingeniería de Sistemas — UPeU
%=============================================================

\documentclass[12pt, a4paper]{article}

\PassOptionsToPackage{table, xcdraw}{xcolor}
\usepackage[utf8]{inputenc}
\usepackage[T1]{fontenc}
\usepackage[spanish, es-nodecimaldot]{babel}
\usepackage{csquotes}
\usepackage{lmodern}
\usepackage{microtype}
\usepackage{setspace}
\onehalfspacing
\usepackage[top=3cm, bottom=2.5cm, left=3cm, right=2.5cm]{geometry}
\usepackage{amsmath}
\usepackage{amssymb}
\usepackage{graphicx}
\usepackage{float}
\usepackage{caption}
\usepackage{subcaption}
\usepackage{tikz}
\usetikzlibrary{shapes.geometric, arrows.meta, positioning,
                fit, backgrounds, calc}
\usepackage{booktabs}
\usepackage{tabularx}
\usepackage{array}
\usepackage{multirow}
\usepackage{longtable}
\renewcommand{\arraystretch}{1.35}
\usepackage{enumitem}
\setlist[itemize]{leftmargin=1.5em, itemsep=2pt}
\setlist[enumerate]{leftmargin=1.5em, itemsep=2pt}
\usepackage{xcolor}
\definecolor{azulPrincipal}{RGB}{0, 70, 127}
\definecolor{azulMedio}{RGB}{0, 102, 179}
\definecolor{azulClaro}{RGB}{204, 229, 255}
\definecolor{grisClaro}{RGB}{245, 245, 245}
\definecolor{verdeOsc}{RGB}{0, 100, 60}
\definecolor{naranjaAlerta}{RGB}{255, 140, 0}
\usepackage{tcolorbox}
\tcbuselibrary{skins, breakable, theorems}

\newtcolorbox{hallazgo}[1][]{
  enhanced,
  colback=verdeOsc!8,
  colframe=verdeOsc,
  fonttitle=\bfseries\color{white},
  coltitle=white,
  colbacktitle=verdeOsc,
  title={Hallazgo},
  #1
}

\usepackage{fancyhdr}
\setlength{\headheight}{15pt}
\pagestyle{fancy}
\fancyhf{}
\fancyhead[L]{\small\textcolor{azulPrincipal}{\textit{Detecci\'{o}n de personas: oclusi\'{o}n y baja visibilidad en videovigilancia}}}
\fancyhead[R]{\small\textcolor{azulPrincipal}{\textit{Art\'{i}culo de Revisi\'{o}n}}}
\fancyfoot[C]{\small\thepage}
\renewcommand{\headrulewidth}{0.4pt}

\usepackage{titlesec}
\titleformat{\section}{\large\bfseries\color{azulPrincipal}}{\thesection.}{0.6em}{}[\titlerule]
\titleformat{\subsection}{\normalsize\bfseries\color{azulMedio}}{\thesubsection.}{0.5em}{}
\titleformat{\subsubsection}{\normalsize\itshape\color{azulPrincipal}}{\thesubsubsection.}{0.4em}{}

\usepackage{hyperref}
\hypersetup{
  colorlinks=true, linkcolor=azulPrincipal,
  citecolor=azulMedio, urlcolor=azulMedio,
}

\usepackage[backend=biber, style=ieee, sorting=nyt, maxbibnames=6]{biblatex}
\addbibresource{referencias.bib}

\raggedbottom


%=============================================================
\begin{document}
%=============================================================

%-------------------------------------------------------------
% PORTADA
%-------------------------------------------------------------
{
\begin{center}
  {\LARGE\bfseries\color{azulPrincipal}
    Modelos para la detecci\'{o}n de personas ante oclusi\'{o}n\\
    y baja visibilidad en sistemas de videovigilancia:\\
    Una revisi\'{o}n sistem\'{a}tica de la literatura}\\[1em]

  {\large
    Celia Patricia Apaza Hilasaca\textsuperscript{1},\quad
    Brayan Raul Condori Quispe\textsuperscript{2},\quad
    Jhan Logan Ramos Quispe\textsuperscript{3}}\\[0.5em]

  {\small
    \textsuperscript{1}Escuela Profesional de Ingenier\'{i}a de Sistemas,
    Universidad Peruana Uni\'{o}n, Juliaca, Per\'{u}.
    \href{mailto:celia.apaza@upeu.edu.pe}{celia.apaza@upeu.edu.pe}\\
    \textsuperscript{2}Escuela Profesional de Ingenier\'{i}a de Sistemas,
    Universidad Peruana Uni\'{o}n, Juliaca, Per\'{u}.
    \href{mailto:brayan.condori@upeu.edu.pe}{brayan.condori@upeu.edu.pe}\\
    \textsuperscript{3}Escuela Profesional de Ingenier\'{i}a de Sistemas,
    Universidad Peruana Uni\'{o}n, Juliaca, Per\'{u}.
    \href{mailto:jhan.ramos@upeu.edu.pe}{jhan.ramos@upeu.edu.pe}}\\[0.8em]

  {\small\textbf{Recibido:} ---/---/---- \quad
         \textbf{Aceptado:} ---/---/---- \quad
         \textbf{Publicado:} ---/---/----}
\end{center}
}

\noindent\rule{\linewidth}{1pt}

%-------------------------------------------------------------
% RESUMEN
%-------------------------------------------------------------
\section*{Resumen}
\noindent
\textbf{Antecedentes:} El continuo aumento del uso de cámaras en la
vigilancia de entornos urbanos e industriales ha generado una
demanda constante y sostenida de algoritmos capaces de localizar
personas con alta precisión bajo condiciones que degradan la calidad
visual en las imágenes provenientes de las cámaras, como la oclusión entre individuos u
objetos y los escenarios con poca iluminación o condiciones
meteorológicas que son inusuales y que a la vez dificultan su lectura e interpretación por la tecnología usada.
\textbf{Objetivo:} Sintetizar los modelos y técnicas documentados
en la literatura científica actual orientados a la detección
automática de personas bajo situaciones de obstrucción visual y
baja visibilidad dentro de plataformas de videovigilancia existentes en su área, e identificar
los conjuntos de datos así como las métricas, limitaciones y vacíos de
conocimiento predominantes en este campo de la tecnología.
\textbf{Metodología:} Se ejecutó una Revisión Sistemática de
Literatura siguiendo el protocolo de Kitchenham y
Charters, con búsquedas en cuatro bases de
datos indexadas: ACM Digital Library, Scopus, ScienceDirect y Web of
Science, acotadas al período 2022--2026. La selección de estos
estudios se realizó en cuatro pasos aplicando criterios de
inclusión y exclusión predefinidos previamente.
\textbf{Resultados:} De 1\,107 registros recuperados, se retuvieron
35 estudios primarios. Los enfoques visiblemente más frecuentes son las
arquitecturas de detección basadas en fusión multiespectral
(RGB + infrarrojo o térmico), las redes con módulos de
atención y las variantes compactas de YOLO optimizadas para trabajar en tiempo
real. Las métricas predominantes son la precisión media
(\textit{mAP}), la tasa de error de detección y los fotogramas
por segundo (\textit{FPS}).
\textbf{Conclusiones:} La fusión multimodal y los mecanismos de
atención son las estrategias con más respaldo empírico 
para gestionar la oclusión y la baja iluminación; sin embargo, aún
persisten carencias en la generalización cuando se trata de escenarios reales y en
la disponibilidad de conjuntos de datos que combinen ambas condiciones
adversas simultáneamente en el mismo área.

\vspace{0.4em}
\noindent\textbf{Palabras clave:} detección de personas,
videovigilancia, oclusión, baja visibilidad, aprendizaje profundo,
redes neuronales convolucionales, detección de peatones.

\vspace{0.4em}\noindent\rule{\linewidth}{0.4pt}

%-------------------------------------------------------------
% ABSTRACT
%-------------------------------------------------------------
\section*{Abstract}
\noindent
\textbf{Background:} The continuous increase in the use of cameras for
surveillance in urban and industrial environments has generated a
constant and sustained demand for algorithms capable of locating
people with high precision under conditions that degrade the visual quality
of camera images, such as occlusion between individuals or
objects, and scenarios with low illumination or unusual weather conditions
that simultaneously hinder their reading and interpretation by the technology used.
\textbf{Objective:} To synthesize the models and techniques documented
in the current scientific literature oriented toward the automatic detection
of people under situations of visual obstruction and low visibility
within video surveillance platforms existing in their field, and to identify
the datasets as well as the metrics, limitations, and predominant knowledge
gaps in this technological field.
\textbf{Methodology:} A Systematic Literature Review was executed following
the Kitchenham and Charters protocol, with searches in four indexed
databases: ACM Digital Library, Scopus, ScienceDirect, and Web of
Science, limited to the 2022--2026 period. The selection of these
studies was carried out in four steps applying previously predefined inclusion
and exclusion criteria.
\textbf{Results:} From 1,107 retrieved records, 35 primary studies
were retained. The visibly most frequent approaches are detection
architectures based on multispectral fusion (RGB + infrared or thermal),
networks with attention modules, and compact variants of YOLO optimized to work in
real time. The predominant metrics are mean average precision
(\textit{mAP}), the detection miss rate, and frames per
second (\textit{FPS}).
\textbf{Conclusions:} Multimodal fusion and attention mechanisms are
the strategies with the most empirical support for managing occlusion
and low illumination; however, shortcomings still persist in generalization
when dealing with real-world scenarios and in the availability of datasets that
combine both adverse conditions simultaneously in the same area.

\vspace{0.4em}
\noindent\textbf{Keywords:} person detection, video surveillance,
occlusion, low visibility, deep learning, convolutional neural
networks, pedestrian detection.

\noindent\rule{\linewidth}{1pt}
\vspace{1em}

%=============================================================
\section{Introducci\'{o}n}
%=============================================================

La vigilancia autom\'{a}tica a traves de c\'{a}maras significa hoy un
componente central en la gesti\'{o}n de la seguridad de las personas, el
control de accesos y el monitoreo de infraestructuras
cr\'{i}ticas~\cite{E24, E32}. La identificaci\'{o}n de personas en
tiempo real dentro de estos sistemas existentes es el requisito m\'{a}s exigente,
ya que condiciona la utilidad de cualquier plataforma de software inteligente de
alerta temprana~\cite{E12, E09}.

Sin embargo, el desempe\~{n}o de los detectores modernos se ve
comprometido cuando las condiciones de captura se alejan de lo
\'{o}ptimo e ideal. La oclusi\'{o}n ---entendida como la obstrucci\'{o}n
parcial o total de un individuo por otro objeto o persona--- y la baja
visibilidad ---propisiada por escasa iluminaci\'{o}n, niebla, lluvia o
condiciones meteorol\'{o}gicas que desfavorecen el trabajo de la tecnologia--- son los factores m\'{a}s
citados en la literatura como fuentes de degradaci\'{o}n del
rendimiento de los detectores~\cite{Dollar2012, E27, E11}. Aunque
arquitecturas como YOLO~\cite{Redmon2016} y Faster
R-CNN~\cite{Ren2015} han alcanzado resultados destacados en condiciones
controladas, la evidencia emp\'{i}rica muestra que su efectividad
disminuye sensiblemente cuando esta ante estos escenarios adversos.

La comunidad cient\'{i}fica ha explorado diversas respuestas ante esta
problem\'{a}tica: la incorporaci\'{o}n de canales espectrales
adicionales ---im\'{a}genes infrarrojas  o t\'{e}rmicas--- para reducir
la dependencia de la luz visible~\cite{E35, E31, E19, E26}; el
dise\~{n}o de m\'{o}dulos de atenci\'{o}n que antepongan las regiones
parcialmente visibles~\cite{E25, E32, E23}; y el desarrollo de
arquitecturas compactas orientadas a la inferencia en tiempo real sobre
hardware con recursos limitados~\cite{E21, E23}. Sin embargo, no
existe a la fecha actual una s\'{i}ntesis estructurada de estos avances
espec\'{i}ficamente acotada al contexto de videovigilancia con
oclusi\'{o}n y baja visibilidad que trabajen de forma simult\'{a}nea.

Esta Revisi\'{o}n Sistem\'{a}tica de Literatura realiza las siguientes
contribuciones: (i) identifica y clasifica los modelos de detecci\'{o}n
de individuos publicados entre el a\~{n}o 2022 y el a\~{n}o 2026 bajo condiciones considerables de
oclusi\'{o}n y baja visibilidad; (ii) sistematiza los conjuntos de
datos, m\'{e}tricas y resultados reportados; y (iii) mapea las
limitaciones y vacios de conocimiento que orientan la investigaci\'{o}n
futura.

El resto del documento se organiza de la siguiente manera: la
Secci\'{o}n~\ref{sec:marco} establece los conceptos fundamentales; la
Secci\'{o}n~\ref{sec:metodo} describe el proceso metodol\'{o}gico; la
Secci\'{o}n~\ref{sec:resultados} presenta los resultados; la
Secci\'{o}n~\ref{sec:discusion} ofrece la discusi\'{o}n; y la
Secci\'{o}n~\ref{sec:conclusiones} cierra con las conclusiones.

\subsection*{Preguntas de Investigaci\'{o}n}

\textbf{Preguntas de investigaci\'{o}n (RQ):}
\begin{enumerate}
  \item \textbf{RQ1:} \textquestiondown Qu\'{e} modelos de detecci\'{o}n
        de personas han sido propuestos para condiciones de oclusi\'{o}n
        y baja visibilidad en sistemas de videovigilancia?
  \item \textbf{RQ2:} \textquestiondown Qu\'{e} t\'{e}cnicas o enfoques
        se utilizan para manejar la oclusi\'{o}n parcial y total en la
        detecci\'{o}n de personas dentro de sistemas de videovigilancia?
  \item \textbf{RQ3:} \textquestiondown Qu\'{e} m\'{e}tricas son
        utilizadas para evaluar el rendimiento de los modelos de
        detecci\'{o}n de personas en condiciones de oclusi\'{o}n y baja
        visibilidad?
  \item \textbf{RQ4:} \textquestiondown Qu\'{e} conjuntos de datos han
        sido utilizados para entrenar y evaluar modelos de detecci\'{o}n
        de personas en condiciones de oclusi\'{o}n y baja visibilidad en
        entornos de videovigilancia?
  \item \textbf{RQ5:} \textquestiondown Cu\'{a}les son las limitaciones
        y desaf\'{i}os reportados por los autores en los modelos de
        detecci\'{o}n de personas ante oclusi\'{o}n y baja visibilidad?
  \item \textbf{RQ6:} \textquestiondown Qu\'{e} vac\'{i}os de
        conocimiento y trabajos futuros han identificado los autores en
        el \'{a}mbito de la detecci\'{o}n de personas en sistemas de
        videovigilancia?
\end{enumerate}

%=============================================================
\section{Marco Conceptual}
\label{sec:marco}
%=============================================================

\subsection{Sistemas de Videovigilancia}

Un sistema de videovigilancia integra dispositivos de captura de
imagen o video, infraestructura de transmisi\'{o}n de datos y
m\'{o}dulos de an\'{a}lisis computacional orientados a la
detecci\'{o}n autom\'{a}tica de events relevantes en el \'{a}rea
monitoreada~\cite{E24, E32}. Su evoluci\'{o}n hacia arquitecturas de
procesamiento en tiempo real, impulsada por el aprendizaje profundo, ha
ampliado su aplicabilidad desde la seguridad perimetral hasta la
gesti\'{o}n del flujo peatonal en espacios p\'{u}blicos~\cite{E12,
E09}. Indicadores como la velocidad de inferencia, la precisi\'{o}n de
detecci\'{o}n y la robustez ante condiciones ambientales variables son
los factores que determinan la utilidad pr\'{a}ctica de estos
sistemas~\cite{Brunetti2018}.

\subsection{Detecci\'{o}n de Personas}

La localizaci\'{o}n autom\'{a}tica de individuos en im\'{a}genes o
secuencias de video constituye una tarea de detecci\'{o}n de objetos
cuyo resultado es un conjunto de cuadros delimitadores que encierran a
cada persona identificada~\cite{Dollar2012}. Los detectores modernos se
agrupan en dos familias: los de dos etapas, que primero generan
propuestas de regi\'{o}n y luego las clasifican ---siendo Faster
R-CNN~\cite{Ren2015} el representative m\'{a}s citado---; y los de una
etapa, que realizan ambas operaciones de forma conjunta para lograr
mayor velocidad, con YOLO~\cite{Redmon2016} como exponente principal.
La elecci\'{o}n entre estas familias involucra un compromiso entre
precisi\'{o}n y eficiencia computacional que resulta cr\'{i}tico en
aplicaciones de vigilancia en tiempo real~\cite{E25, E23}.

\subsection{Oclusi\'{o}n en Videovigilancia}

La oclusi\'{o}n aparece cuando un individuo queda parcial o totalmente
encubierto por otro objeto, persona o elemento arquitect\'{o}nico del
entorno que est\'{a} en el mismo \'{a}rea~\cite{E30}. Se distinguen
dos formas: la oclusi\'{o}n parcial, en la que solo una parte
del cuerpo permanece visible desde la perspectiva de la camara ---situaci\'{o}n habitual en pasillos
concurridos o multitudes de gente~\cite{E28, E29}---; y la oclusi\'{o}n
total, en la que el individuo desaparece completamente del plano visual
de la c\'{a}mara, cuyo caso representa el mayor reto para los
algoritmos que existen en la actualidad~\cite{E27, E11}.

\subsection{Baja Visibilidad en Videovigilancia }

El t\'{e}rmino baja visibilidad abarca las condiciones del entorno
que reducen la confiabilidad de la se\~{n}al de la imagen captada por el
sensor. Entre sus causas se encuentran la escasez de iluminaci\'{o}n
ambiental nocturna, la presencia de aerosoles como niebla o humo, y
las precipitaciones intensas~\cite{E21, E12, E13}. Para afrontar estas
limitaciones se han explorado dos l\'{i}neas complementarias: el empleo
de sensores t\'{e}rmicos o infrarrojos, que generan se\~{n}ales
robustas ante cambios de iluminaci\'{o}n~\cite{E34, E35, E14}; y las
t\'{e}cnicas de mejora de imagen que incrementan el contraste antes de
aplicar el detector~\cite{E21}.

%=============================================================
\section{Proceso Metodol\'{o}gico}
\label{sec:metodo}
%=============================================================

Esta revisi\'{o}n sigue el protocolo de~\cite{Kitchenham2007} para la
realizaci\'{o}n de revisiones sistem\'{a}ticas en ingenier\'{i}a de
software, adaptado al \'{a}rea de visi\'{o}n por computadora aplicada
a videovigilancia, y sigue los lineamientos de reporte PRISMA
2020~\cite{Page2021}.

\subsection{Protocolo de Revisi\'{o}n (PICOC)}

\begin{table}[H]
\centering
\caption{Marco PICOC de la revisi\'{o}n sistem\'{a}tica.}
\label{tab:picoc}
\begin{tabularx}{\linewidth}{>{\bfseries}lX}
\toprule
Elemento & Descripci\'{o}n \\
\midrule
Poblaci\'{o}n    & Sistemas de detecci\'{o}n e indentificacion de personas, sistemas usados de
                   videovigilancia, sistemas de detecci\'{o}n de peatones. \\
Intervenci\'{o}n & Modelos de detecci\'{o}n, modelos de aprendizaje
                   profundo, algoritmos de visi\'{o}n por computadora,
                   t\'{e}cnicas de detecci\'{o}n de objetos, m\'{e}todos
                   de manejo de oclusi\'{o}n, m\'{e}todos de
                   detecci\'{o}n en baja visibilidad. \\
Comparaci\'{o}n  & No aplica. \\
Resultado        & Precisi\'{o}n de detecci\'{o}n, \textit{precision},
                   \textit{recall}, F1-score, rendimiento del modelo,
                   robustez, detecci\'{o}n en tiempo real. \\
Contexto         & Entornos de videovigilancia, sistemas de monitoreo
                   de seguridad, escenarios al aire libre y en interiores. \\
\bottomrule
\end{tabularx}
\end{table}

\subsection{Bases de Datos y Fuentes}

\begin{itemize}
  \item ACM Digital Library (\url{https://dl.acm.org})
  \item Scopus (\url{https://www.scopus.com})
  \item Web of Science (\url{https://www.webofscience.com})
  \item ScienceDirect/Elsevier (\url{https://www.sciencedirect.com})
\end{itemize}

\subsection{Cadena de B\'{u}squeda}

La cadena de b\'{u}squeda se construy\'{o} a partiendo de los
t\'{e}rminos PICOC y sus sin\'{o}nimos, conectados con operadores
booleanos.

\noindent\textbf{Cadena base (formato general):}
\begin{center}
\fbox{\parbox{0.95\linewidth}{\small\ttfamily\raggedright
(``person detection'' OR ``pedestrian detection'' OR ``human detection'')
AND (``video surveillance'' OR ``surveillance system'' OR ``CCTV'' OR
``security camera'')
AND (``occlusion'' OR ``partial occlusion'' OR ``full occlusion'')
AND (``low visibility'' OR ``low illumination'' OR ``nighttime'' OR
``fog'' OR ``rain'' OR ``darkness'')
AND (``detection model'' OR ``detection method'' OR ``detection technique''
OR ``deep learning'' OR ``neural network'')
}}
\end{center}

En la Tabla~\ref{tab:busqueda} se presentan los resultados de la
b\'{u}squeda por base de datos y las cadenas espec\'{i}ficas empleadas.

\begin{table}[H]
\centering
\caption{Resultados de la b\'{u}squeda por base de datos.}
\label{tab:busqueda}
\small
\begin{tabularx}{\linewidth}{lXc}
\toprule
\textbf{Base de datos} & \textbf{Cadena espec\'{i}fica} & \textbf{Total} \\
\midrule
ACM Digital Library &
{\ttfamily\raggedright [[Title: ``person detection''] OR [Title: ``pedestrian
detection''] OR [Title: ``human detection'']] AND [[Abstract:
``surveillance''] OR [Abstract: ``occlusion''] OR [Abstract: ``low
visibility''] OR [Abstract: ``infrared''] OR [Abstract: ``nighttime'']]}
AND Publication Date: (01/01/2022 TO 12/31/2026)
& 151 \\
\midrule
Scopus &
\small TITLE-ABS-KEY((``person detection'' OR ``pedestrian detection'' OR
``human detection'' OR ``people detection'') AND (``video surveillance''
OR ``surveillance'' OR ``CCTV'' OR ``security camera'' OR ``monitoring'')
AND (``occlusion'' OR ``low visibility'' OR ``low illumination'' OR
``nighttime'' OR ``fog'' OR ``rain'') AND (``deep learning'' OR ``neural
network'' OR ``convolutional'' OR ``object detection'' OR ``feature
extraction'')) AND PUBYEAR $>$ 2021 AND DOCTYPE IN (``ar'',``cp'') AND
LANGUAGE IN (``English'',``Spanish'')
& 162 \\
\midrule
ScienceDirect &
(``pedestrian detection'' OR ``person detection'') AND (``video
surveillance'' OR ``CCTV'') AND (``occlusion'' OR ``low visibility'')
AND (``deep learning'' OR ``neural network'')
& 492 \\
\midrule
Web of Science &
TS=((``person detection'' OR ``pedestrian detection'' OR ``human
detection'') AND (``video surveillance'' OR ``CCTV'') AND (``occlusion''
OR ``low visibility'' OR ``low illumination'') AND (``deep learning'' OR
``neural network'') AND (``real-time'' OR ``accuracy'' OR
``performance'')) AND PY=(2022--2026)
& 302 \\
\midrule
\textbf{TOTAL} & & \textbf{1\,107} \\
\bottomrule
\end{tabularx}
\end{table}

\subsection{Criterios de Inclusi\'{o}n y Exclusi\'{o}n}

\begin{table}[H]
\centering
\caption{Criterios de selecci\'{o}n de estudios primarios.}
\label{tab:criterios}
\begin{tabularx}{\linewidth}{lX}
\toprule
\textbf{C\'{o}digo} & \textbf{Criterio} \\
\midrule
\multicolumn{2}{l}{\cellcolor{azulClaro!40}\textbf{Inclusi\'{o}n}} \\
CI-1 & Se aceptar\'{a}n art\'{i}culos que propongan o eval\'{u}en modelos
       de detecci\'{o}n de personas bajo condiciones de oclusi\'{o}n y/o
       baja visibilidad en sistemas de videovigilancia. \\
CI-2 & Se considerar\'{a}n todos los art\'{i}culos publicados entre
       2022 y 2026. \\
CI-3 & Se consideran todos los art\'{i}culos que provenienen de bases de
       datos digitales indexadas. \\
CI-4 & Los art\'{i}culos deben pertenecer al \'{a}rea de visi\'{o}n por
       computadora, aprendizaje profundo o sistemas de videovigilancia. \\
CI-5 & Se aceptar\'{a}n art\'{i}culos que provengan de revistas
       cient\'{i}ficas y conferencias indexadas. \\
\midrule
\multicolumn{2}{l}{\cellcolor{red!10}\textbf{Exclusi\'{o}n}} \\
CE-1 & Van a ser excluidos los art\'{i}culos duplicados. \\
CE-2 & Van a ser excluidos los art\'{i}culos cuyo t\'{i}tulo no tenga
       relaci\'{o}n con el objeto de estudio de la revisi\'{o}n. \\
CE-3 & Van a ser excluidos los estudios secundarios, estudios terciarios,
       revisiones narrativas sin evaluaci\'{o}n emp\'{i}rica y
       res\'{u}menes. \\
CE-4 & Van a ser rechazados los art\'{i}culos que no se encuentren en
       idioma ingl\'{e}s o espa\~{n}ol. \\
CE-5 & Van a ser rechazados los art\'{i}culos de contenido similar,
       qued\'{a}ndose unicamente con los que tengan el contenido m\'{a}s
       completo. \\
\bottomrule
\end{tabularx}
\end{table}

\subsection{Proceso de Selecci\'{o}n}

El proceso de selecci\'{o}n se realiz\'{o} siguiendo las tres fases
del protocolo propuesto por Kitchenham~\cite{Kitchenham2007}:

\begin{enumerate}
  \item \textbf{Fase 1 --- Planificaci\'{o}n:} identificaci\'{o}n de
        la considerable necesidad de la revisi\'{o}n, elaboraci\'{o}n y
        evaluaci\'{o}n del protocolo.
  \item \textbf{Fase 2 --- Ejecuci\'{o}n:} identificaci\'{o}n de
        estudios, selecci\'{o}n, evaluaci\'{o}n de la calidad,
        extracci\'{o}n y s\'{i}ntesis de datos.
  \item \textbf{Fase 3 --- Documentaci\'{o}n:} redacci\'{o}n,
        validaci\'{o}n y publicaci\'{o}n del informe de revisi\'{o}n.
\end{enumerate}

\begin{figure}[H]
\centering
\begin{tikzpicture}[
  fase/.style={draw=azulPrincipal, fill=azulPrincipal, text=white,
               font=\bfseries\small, text width=3.2cm, align=center,
               minimum height=0.9cm, rounded corners=3pt},
  act/.style={draw=azulPrincipal, fill=azulClaro!40, font=\small,
              text width=3.2cm, align=center, minimum height=0.75cm,
              rounded corners=2pt},
  arr/.style={-Stealth, thick, color=azulPrincipal}
]
  \node[fase] (f1)        {FASE 1\\Planificaci\'{o}n};
  \node[act,  below=0.3cm of f1]  (p1)    {Necesidad de la revisi\'{o}n};
  \node[act,  below=0.2cm of p1]  (p2)    {Elaborar protocolo};
  \node[act,  below=0.2cm of p2]  (p3)    {Evaluar protocolo};
  \node[fase, below=0.55cm of p3] (f2)    {FASE 2\\Ejecuci\'{o}n};
  \node[act,  below=0.3cm of f2]  (e1)    {Identificar estudios (1107)};
  \node[act,  below=0.2cm of e1]  (e2)    {Selecci\'{o}n: 4 pasos (1107$\to$35)};
  \node[act,  below=0.2cm of e2]  (e3)    {Evaluar calidad (CQI)};
  \node[act,  below=0.2cm of e3]  (e4)    {Extraer datos};
  \node[act,  below=0.2cm of e4]  (e5)    {Sintetizar datos};
  \node[fase, below=0.55cm of e5] (f3)    {FASE 3\\Documentaci\'{o}n};
  \node[act,  below=0.3cm of f3]  (d1)    {Redactar informe RSL};
  \node[act,  below=0.2cm of d1]  (d2)    {Validar informe};
  \node[act,  below=0.2cm of d2]  (d3)    {Publicar informe};
  \foreach \a/\b in {f1/p1, p1/p2, p2/p3, p3/f2,
                     f2/e1, e1/e2, e2/e3, e3/e4, e4/e5, e5/f3,
                     f3/d1, d1/d2, d2/d3}
    \draw[arr] (\a) -- (\b);
\end{tikzpicture}
\caption{Proceso de Revisión Sistemática de Literatura seg\'{u}n~\cite{Kitchenham2007}.}
\label{fig:prisma}
\end{figure}

La Tabla~\ref{tab:pasos} muestra los criterios aplicados en cada paso
y la Tabla~\ref{tab:seleccion} presenta los resultados cuantitativos
por base de datos.

\begin{table}[H]
\centering
\caption{Procedimientos y criterios de selecci\'{o}n.}
\label{tab:pasos}
\begin{tabularx}{\linewidth}{lX}
\toprule
\textbf{Procedimiento} & \textbf{Criterio de selecci\'{o}n} \\
\midrule
Paso 1 & CI-3, CI-4, CI-5, CE-1, CE-3 \\
Paso 2 & CE-2, CE-3, CE-4 \\
Paso 3 & CE-3, CE-5 \\
Paso 4 & CI-1, CI-2 \\
\bottomrule
\end{tabularx}
\end{table}

\begin{table}[H]
\centering
\caption{Resultados del proceso de selecci\'{o}n de estudios.}
\label{tab:seleccion}
\begin{tabularx}{\linewidth}{Xccccc}
\toprule
\textbf{Base de datos} & \textbf{Inicial} &
\textbf{P1} & \textbf{P2} & \textbf{P3} & \textbf{P4} \\
\midrule
ACM Digital Library & 151 & 42 & 19 & 14 &  8 \\
Scopus              & 162 & 115 & 104 & 103 & 12 \\
ScienceDirect       & 492 & 344 & 198 & 198 & 12 \\
Web of Science      & 302 & 181 & 158 & 158 &  3 \\
\midrule
\textbf{Total} & \textbf{1\,107} & \textbf{682} &
\textbf{479} & \textbf{473} & \textbf{35} \\
\bottomrule
\end{tabularx}
\end{table}

\subsection{Evaluaci\'{o}n de Calidad}

Cada estudio primario se evalu\'{o} mediante la lista de
verificaci\'{o}n de calidad adaptada a la escala de
Rouhani~\cite{Kitchenham2013}, con escala: S\'{i}~=~1 / Parcialmente
=~0{,}5 / No~=~0. La Tabla~\ref{tab:calidad} detalla los \'{i}tems
considerados.

\begin{table}[H]
\centering
\caption{Instrumento de evaluaci\'{o}n de calidad (CQI).}
\label{tab:calidad}
\begin{tabularx}{\linewidth}{c>{\raggedright\arraybackslash}Xl}
\toprule
\textbf{\'{I}tem} & \textbf{Criterio} & \textbf{Escala} \\
\midrule
CQ-1 & \textquestiondown El m\'{e}todo o modelo de detecci\'{o}n de
       personas ha sido documentado apropiadamente?
     & S\'{i} (1) / Parcial (0{,}5) / No (0) \\
CQ-2 & \textquestiondown El estudio eval\'{u}a el modelo en condiciones
       de oclusi\'{o}n y/o baja visibilidad de manera expl\'{i}cita?
     & S\'{i} (1) / Parcial (0{,}5) / No (0) \\
CQ-3 & \textquestiondown Se han documentado las limitaciones del
       estudio de manera clara?
     & S\'{i} (1) / Parcial (0{,}5) / No (0) \\
CQ-4 & \textquestiondown Los aportes del estudio para la comunidad
       cient\'{i}fica, acad\'{e}mica o para la industria han sido
       descritos?
     & S\'{i} (1) / Parcial (0{,}5) / No (0) \\
CQ-5 & \textquestiondown Los resultados han contribuido a responder las
       preguntas de investigaci\'{o}n planteadas?
     & S\'{i} (1) / Parcial (0{,}5) / No (0) \\
\midrule
\multicolumn{2}{l}{\textbf{Puntuaci\'{o}n m\'{i}nima de inclusi\'{o}n}}
     & $\geq 2{,}5$ / 5{,}0 \\
\bottomrule
\end{tabularx}
\end{table}

\subsection{Extracci\'{o}n de Datos}

Los datos de cada estudio primario se extrajeron en el formulario de
la Tabla~\ref{tab:extraccion}:

\begin{table}[H]
\centering
\caption{Formulario de extracci\'{o}n de datos.}
\label{tab:extraccion}
\begin{tabularx}{\linewidth}{lX}
\toprule
\textbf{Campo} & \textbf{Descripci\'{o}n} \\
\midrule
ID                       & Identificador \'{u}nico (E01\ldots E35) \\
T\'{i}tulo del art\'{i}culo & T\'{i}tulo completo del estudio \\
Autor(es)                & Apellidos y nombres de los autores \\
A\~{n}o de publicaci\'{o}n & A\~{n}o de publicaci\'{o}n del
                             art\'{i}culo \\
Fuente de publicaci\'{o}n & Revista o conferencia de publicaci\'{o}n \\
Tipo de publicaci\'{o}n  & Art\'{i}culo de revista / Art\'{i}culo de
                           conferencia / Revisi\'{o}n \\
Base de datos de origen  & ACM / Scopus / ScienceDirect /
                           Web of Science \\
Modelo/t\'{e}cnica       & Denominaci\'{o}n y familia del detector \\
Tipo de oclusi\'{o}n     & Parcial / Total / Ambas / No especifica \\
Condici\'{o}n de visibilidad & Baja iluminaci\'{o}n / Niebla / Lluvia
                               / Nocturno / Varias / No especifica \\
M\'{e}tricas utilizadas  & mAP, miss rate, F1-score, FPS, etc. \\
Conjunto de datos        & Benchmark(s) empleado(s) \\
Resultados principales   & Valores obtenidos en las m\'{e}tricas \\
Limitaciones reportadas  & Restricciones declaradas por los autores \\
Vac\'{i}os y trabajo futuro & Brechas identificadas y propuestas \\
Dominio de aplicaci\'{o}n & Interior / Exterior / Ambos \\
Calidad del estudio (CQI) & Puntuaci\'{o}n sobre 5{,}0 \\
\bottomrule
\end{tabularx}
\end{table}

%=============================================================
\section{Resultados}
\label{sec:resultados}
%=============================================================

\subsection{Descripci\'{o}n General de los Estudios}

El proceso de la selecci\'{o}n redujo los 1\,107 registros iniciales a
\textbf{35 estudios primarios} cubriendo el per\'{i}odo 2022--2026.
ACM Digital Library aport\'{o} 8 estudios (22{,}9\%), Scopus 12
(34{,}3\%), ScienceDirect 12 (34{,}3\%) y Web of Science 3 (8{,}6\%).
A continuaci\'{o}n, la Tabla~\ref{tab:estudios} resume de forma extendida los estudios seleccionados.

\begin{singlespace}
\small
\begin{longtable}{clp{7.5cm}ccl}
\caption{Resumen de estudios primarios seleccionados.} \label{tab:estudios} \\
\toprule
\textbf{ID} & \textbf{Ref.} & \textbf{Enfoque principal} & \textbf{A\~{n}o} & \textbf{CQI} & \textbf{RQ} \\
\midrule
\endfirsthead
\caption[]{Resumen de estudios primarios seleccionados (continuaci\'{o}n).} \\
\toprule
\textbf{ID} & \textbf{Ref.} & \textbf{Enfoque principal} & \textbf{A\~{n}o} & \textbf{CQI} & \textbf{RQ} \\
\midrule
\endhead
\midrule
\endfoot
\bottomrule
\endlastfoot
E01 & \cite{E01} & Ataque adversarial infrarrojo nocturno (ACM) & 2023 & 4.0 & RQ1,RQ2 \\
E02 & \cite{E02} & Fusi\'{o}n sem\'{a}ntica visible-infrarrojo (ACM) & 2025 & 4.5 & RQ1,RQ2,RQ4 \\
E03 & \cite{E03} & Detecci\'{o}n peatonal oclusi\'{o}n cruzada (ACM) & 2024 & 4.0 & RQ1,RQ2 \\
E04 & \cite{E04} & Fusi\'{o}n t\'{e}rmica-RGB clima adverso (ACM) & 2023 & 4.0 & RQ1,RQ2,RQ4 \\
E05 & \cite{E05} & Atenci\'{o}n oclusi\'{o}n multitud vigilancia (ACM) & 2022 & 4.0 & RQ1,RQ2 \\
E06 & \cite{E06} & Peatones tiempo real baja luz CCTV (ACM) & 2024 & 4.5 & RQ1,RQ3,RQ4 \\
E07 & \cite{E07} & Alineaci\'{o}n modal visible-infrarrojo (ACM) & 2024 & 4.0 & RQ1,RQ2,RQ4 \\
E08 & \cite{E08} & Mejora infrarrojo vigilancia nocturna (ACM) & 2025 & 4.5 & RQ1,RQ3 \\
E09 & \cite{E09} & Fusi\'{o}n m\'{a}scaras residuales multi-vista & 2026 & 4.5 & RQ1,RQ2 \\
E10 & \cite{E10} & LF-DETR RGB-IR a\'{e}reo & 2026 & 4.5 & RQ1,RQ2,RQ4 \\
E11 & \cite{E11} & Oclusi\'{o}n adaptativa en multitudes & 2026 & 4.5 & RQ1,RQ2,RQ5 \\
E12 & \cite{E12} & NP-YOLO nocturno vigilancia & 2025 & 4.5 & RQ1,RQ3,RQ4 \\
E13 & \cite{E13} & SlowFast mejorado nocturno portuario & 2025 & 4.0 & RQ1,RQ3 \\
E14 & \cite{E14} & Detecci\'{o}n t\'{e}rmica atenci\'{o}n canal-espacial & 2025 & 4.0 & RQ1,RQ2,RQ4 \\
E15 & \cite{E15} & Persona + arma vigilancia inteligente & 2025 & 3.5 & RQ1 \\
E16 & \cite{E16} & Detecci\'{o}n t\'{e}rmica kurtosis + YOLOv8 & 2025 & 4.0 & RQ1,RQ4 \\
E17 & \cite{E17} & Fusi\'{o}n multiespectral diferencial & 2025 & 4.5 & RQ1,RQ2,RQ4 \\
E18 & \cite{E18} & PYOLO oclusi\'{o}n peatonal & 2025 & 4.5 & RQ1,RQ2 \\
E19 & \cite{E19} & Multiespectral CNN fusi\'{o}n canal & 2025 & 4.0 & RQ1,RQ2,RQ4 \\
E20 & \cite{E20} & Transformer recurrente detecci\'{o}n t\'{e}rmica & 2025 & 4.0 & RQ1,RQ4 \\
E21 & \cite{E21} & Mejora imagen + YOLO nocturno & 2026 & 4.5 & RQ1,RQ3 \\
E22 & \cite{E22} & Sem\'{a}ntica visual-ling\"{u}\'{i}stica peatonal & 2025 & 4.0 & RQ1,RQ3 \\
E23 & \cite{E23} & YOLOv8s ligero peque\~{n}a escala & 2025 & 4.5 & RQ1,RQ3,RQ5 \\
E24 & \cite{E24} & SUSAN detecci\'{o}n violencia vigilancia & 2025 & 4.0 & RQ1 \\
E25 & \cite{E25} & Ms-VLPD multiescala peatonal & 2025 & 4.0 & RQ1,RQ3 \\
E26 & \cite{E26} & EAFF-Net fusi\'{o}n dual RGB-IR & 2025 & 4.5 & RQ1,RQ2,RQ4 \\
E27 & \cite{E27} & Esqueleto + IA ca\'{i}das con oclusi\'{o}n & 2025 & 4.0 & RQ1,RQ2,RQ5 \\
E28 & \cite{E28} & DETR adaptativo multitudes & 2025 & 4.0 & RQ1,RQ2 \\
E29 & \cite{E29} & RC-DETR multitudes contrastivo & 2025 & 4.0 & RQ1,RQ2 \\
E30 & \cite{E30} & PPM multi-vista peatonal & 2024 & 4.0 & RQ1,RQ2 \\
E31 & \cite{E31} & Multiespectral atenci\'{o}n cruzada & 2024 & 4.5 & RQ1,RQ2,RQ4 \\
E32 & \cite{E32} & PFEL-Net multiescala ligero & 2024 & 4.0 & RQ1,RQ3 \\
E33 & \cite{E33} & ML video transporte vigilancia & 2026 & 3.5 & RQ1,RQ3 \\
E34 & \cite{E34} & Faster R-CNN + SSD t\'{e}rmico a\'{e}reo & 2022 & 4.5 & RQ1,RQ4 \\
E35 & \cite{E35} & Multiespectral DL una etapa & 2022 & 4.5 & RQ1,RQ2,RQ4 \\
\end{longtable}
\end{singlespace}

\normalsize

\subsection{Resultados por Pregunta de Investigaci\'{o}n}

\subsubsection{RQ1: \textquestiondown Qu\'{e} modelos de detecci\'{o}n
de personas han sido propuestos para condiciones de oclusi\'{o}n y baja
visibilidad en sistemas de videovigilancia?}

El an\'{a}lisis de los 35 estudios primarios permite agrupar los
modelos identificados en cuatro familias principales. Las
\textbf{arquitecturas basadas en YOLO} (YOLOv8s, NP-YOLO, PYOLO)
representan el enfoque de una sola etapa m\'{a}s frecuente, valorado
por su capacidad de procesar m\'{u}ltiples fotogramas por segundo con
recursos computacionales moderados~\cite{E21,E18,E12,E23}. Los
\textbf{detectores de fusi\'{o}n multiespectral} combinan im\'{a}genes
del espectro visible con canales infrarrojos o t\'{e}rmicos, aportando
robustez ante condiciones nocturnas y de niebla; esto ofrece
informaci\'{o}n complementaria independiente de la iluminaci\'{o}n
ambiental~\cite{E35,E31,E19,E17,E07,E04,E26}. Las \textbf{redes con
m\'{o}dulos de atenci\'{o}n} incorporan capas de atenci\'{o}n espacial
o de canal para focalizar la representaci\'{o}n en las regiones m\'{a}s
informativas, incrementando la sensibilidad ante figuras parcialmente
encubiertas~\cite{E25, E05, E14}. Finalmente, los \textbf{enfoques
multi-vista} aprovechan el uso de varias c\'{a}maras simult\'{a}neas
para resolver ambig\"{u}edades originadas por la oclusi\'{o}n mutua
entre personas~\cite{E09, E30, E10}.

\begin{hallazgo}[]
  Los modelos de fusi\'{o}n multiespectral (RGB + IR/t\'{e}rmico) y las
  versiones existentes de YOLO constituyen los enfoques con mayor
  presencia en la literatura revisada, identificados en el 63\,\% de
  los estudios primarios. La fusi\'{o}n multimodal es la estrategia
  predominante para obtener los mejores resultados ante la baja visibilidad
  en horas de la noche, mientras que los m\'{o}dulos de atenci\'{o}n son la
  respuesta principalmente elejida ante la oclusi\'{o}n
  parcial~\cite{E35, E31, E17, E19, E07, E26}.
\end{hallazgo}

\subsubsection{RQ2: \textquestiondown Qu\'{e} t\'{e}cnicas o enfoques
se utilizan para manejar la oclusi\'{o}n parcial y total en la
detecci\'{o}n de personas dentro de sistemas de videovigilancia?}

Los estudios primarios documentan estas cinco estrategias complementarias
para el manejo de la oclusi\'{o}n:
\begin{enumerate}
  \item \textbf{Predicci\'{o}n de partes anatomicas corporales visibles:} el
        detector genera cuadros parciales sobre las regiones expuestas ante la camara
        (cabeza, abdomen) y supone la posici\'{o}n global a partir de
        ellas~\cite{E27}.
  \item \textbf{Estimaci\'{o}n de pose mediante esqueletos:} la
        reconstrucci\'{o}n del esqueleto corporal permite inferir de manera eficar la postura incluso cuando parte del cuerpo
        est\'{a} porcentualmente ocluida~\cite{E27}.
  \item \textbf{Supresi\'{o}n de no m\'{a}ximos adaptada (NMS):}
        variantes que an sido dise\~{n}adas para preservar detecciones de peatones
        solapados en escenas densas que dificultan la vision~\cite{E29, E28, E11}.
  \item \textbf{Fusi\'{o}n de perspectivas m\'{u}ltiples:} el uso
        simult\'{a}neo de una significativa cantidad de c\'{a}maras resuelve las oclusiones que
        aparecen en una sola vista~\cite{E09, E30, E10}.
  \item \textbf{M\'{o}dulos de atenci\'{o}n cruzada:} mecanismos que
        ponderan las regiones parcialmente visibles y disminuyen el
        ruido de fondo~\cite{E31, E17, E14}.
\end{enumerate}

\begin{hallazgo}[]
  La fusi\'{o}n multi-vista y los m\'{o}dulos de atenci\'{o}n cruzada
  son las estrategias con mayor respaldo emp\'{i}rico para escenarios
  de oclusi\'{o}n severa. Los enfoques multi-vista reportan mejoras
  de hasta un 8\,\% en mAP respecto a detectores de referencia en
  escenarios de oclusi\'{o}n densa~\cite{E09, E11}.
\end{hallazgo}

\subsubsection{RQ3: \textquestiondown Qu\'{e} m\'{e}tricas son
utilizadas para evaluar el rendimiento de los modelos de detecci\'{o}n
de personas en condiciones de oclusi\'{o}n y baja visibilidad?}

Las m\'{e}tricas reportadas con mayor frecuencia en los estudios
primarios son:
\begin{itemize}
  \item \textbf{mAP} (\textit{mean Average Precision}): m\'{e}trica de
        referencia en detecci\'{o}n de objetos, empleada en la
        mayor\'{i}a de los estudios~\cite{E31, E25, E10}.
  \item \textbf{Tasa de error logar\'{i}tmica} (\textit{log-average
        miss rate}): indicador espec\'{i}fico del benchmark Caltech,
        utilizado para evaluar el comportamiento ante peatones
        parcialmente visibles~\cite{E29, E28}.
  \item \textit{Precisi\'{o}n y exhaustividad} (\textit{precision} y
        \textit{recall}): m\'{e}tricas complementarias para caracterizar
        el equilibrio entre falsas alarmas y detecciones
        perdidas~\cite{E21, E25, E33}.
  \item \textbf{F1-score:} media arm\'{o}nica de las dos anteriores,
        utilizada cuando el conjunto de datos presenta desequilibrio de
        clases~\cite{E24, E23}.
  \item \textbf{FPS} (\textit{frames per second}): indicador de
        velocidad de inferencia, cr\'{i}tico en aplicaciones de
        vigilancia en tiempo real~\cite{E12, E23}.
\end{itemize}

\begin{hallazgo}[]
  El mAP y el FPS son las m\'{e}tricas prioritarias en los estudios
  revisados: el primero eval\'{u}a la precisi\'{o}n global del
  detector, mientras que el segundo determina su viabilidad para
  aplicaciones de videovigilancia en tiempo real con m\'{u}ltiples
  c\'{a}maras~\cite{E21, E12, E23}.
\end{hallazgo}

\subsubsection{RQ4: \textquestiondown Qu\'{e} conjuntos de datos han
sido utilizados para entrenar y evaluar modelos de detecci\'{o}n de
personas en condiciones de oclusi\'{o}n y baja visibilidad en entornos
de videovigilancia?}

Los benchmarks m\'{a}s frecuentes identificados en los estudios
primarios son:
\begin{itemize}
  \item \textbf{KAIST Multispectral Pedestrian Dataset:} par de
        im\'{a}genes visibles e infrarrojas para escenarios diurnos y
        nocturnos~\cite{E35, E31, E17, E26}.
  \item \textbf{FLIR Thermal Dataset:} im\'{a}genes t\'{e}rmicas para
        detecci\'{o}n con escasa iluminaci\'{o}n~\cite{E34, E14, E16}.
  \item \textbf{Caltech Pedestrian Dataset:} est\'{a}ndar de referencia
        para detecci\'{o}n peatonal urbana con anotaciones de nivel de
        oclusi\'{o}n~\cite{E29, E28, E11}.
  \item \textbf{CityPersons:} derivado de Cityscapes, con peatones en
        entornos urbanos reales y diversas condiciones de
        oclusi\'{o}n~\cite{E22, E25}.
  \item \textbf{COCO:} conjunto de prop\'{o}sito general empleado para
        preentrenamiento y transferencia de dominio~\cite{E20, E32}.
  \item \textbf{MOT Challenge:} para evaluaci\'{o}n de seguimiento y
        detecci\'{o}n multi-persona~\cite{E24}.
\end{itemize}

\begin{hallazgo}[]
  KAIST y FLIR son los datasets m\'{a}s utilizados en estudios de
  detecci\'{o}n bajo baja visibilidad, mientras que Caltech y
  CityPersons predominan en evaluaciones de oclusi\'{o}n. Ninguno de
  los datasets existentes combina simult\'{a}neamente oclusi\'{o}n
  severa y condiciones meteorol\'{o}gicas adversas, lo que representa
  una brecha relevante identificada en m\'{u}ltiples
  estudios~\cite{E31, E17, E07}.
\end{hallazgo}

\subsubsection{RQ5: \textquestiondown Cu\'{a}les son las limitaciones y
desaf\'{i}os reportados por los autores en los modelos de
detecci\'{o}n de personas ante oclusi\'{o}n y baja visibilidad?}

El an\'{a}lisis de las secciones de limitaciones de los estudios
primarios revela cinco categor\'{i}as recurrentes:

\textbf{Dependencia de hardware especializado.} Los modelos que
integran sensores infrarrojos o t\'{e}rmicos requieren equipamiento
adicional que eleva el coste de implementaci\'{o}n~\cite{E14, E26}.
\textbf{Fragilidad ante oclusi\'{o}n considerablemente alta.} Cuando m\'{a}s del
70\,\% de la silueta corporal queda totalmente ocluida, la mayor\'{i}a de los
detectores indican una ca\'{i}da significativa en la tasa de
detecci\'{o}n~\cite{E27, E11}.
\textbf{Brecha entre banco de pruebas con condiciones ideales y entorno real.} Los modelos
optimizados sobre benchmarks controlados no mantienen continuamente sus niveles de
rendimiento habituales al ser desplegados en escenarios con condiciones
variables no presentes en el conjunto de
entrenamiento anteriormente usado~\cite{E33, E23}.
\textbf{Latencia computacional.} Los detectores de mayor precisi\'{o}n
suelen demandar recursos que superan la capacidad de los dispositivos
embebidos, limitando su aplicabilidad en redes de vigilancia
descentralizadas~\cite{E32, E23}.
\textbf{Escasez de datos etiquetados para condiciones extremas.}
La anotaci\'{o}n de im\'{a}genes con oclusi\'{o}n densa, niebla o
lluvia es costosa y los conjuntos existentes no cubren adecuadamente
estas situaciones~\cite{E31, E07}.

\begin{hallazgo}[]
  La brecha entre rendimiento en benchmark y rendimiento en entorno real
  es la limitaci\'{o}n m\'{a}s citada en los estudios revisados
  (presente en el 47\,\% de los art\'{i}culos), seguida por la demanda
  computacional de los modelos m\'{a}s precisos, lo que dificulta su
  despliegue en sistemas de vigilancia con hardware
  limitado~\cite{E27, E33, E11, E23}.
\end{hallazgo}

\subsubsection{RQ6: \textquestiondown Qu\'{e} vac\'{i}os de conocimiento
y trabajos futuros han identificado los autores en el \'{a}mbito de la
detecci\'{o}n de personas en sistemas de videovigilancia?}

Los autores de los estudios primarios se\~{n}alan de manera recurrente
las siguientes \'{a}reas:

\textbf{Datasets que combinen m\'{u}ltiples condiciones adversas.} Se
requieren colecciones de datos que presenten simult\'{a}neamente
oclusi\'{o}n, baja iluminaci\'{o}n y fen\'{o}menos meteorol\'{o}gicos,
para entrenar modelos m\'{a}s robustos~\cite{E31, E07}.
\textbf{Modelos unificados para oclusi\'{o}n y baja visibilidad.} La
mayor\'{i}a de las propuestas actuales abordan cada condici\'{o}n por
separado; existe una brecha en arquitecturas capaces de gestionar ambas
de forma conjunta~\cite{E20, E26}.
\textbf{Adaptaci\'{o}n de dominio.} Los detectores entrenados en un
entorno espec\'{i}fico presentan una caida de eficacia al aplicarse en
contextos distintos; se necesitan estrategias de transferencia considerablemente m\'{a}s
robustas~\cite{E09, E30}.
\textbf{Eficiencia para dispositivos embebidos.} Se requiere el
dise\~{n}o de arquitecturas que mantengan alta precisi\'{o}n con bajo
consumo de hardaware computacional, haciendo posible el procesamiento en el borde de la
red (\textit{edge computing})~\cite{E32, E23}.
\textbf{Evaluaci\'{o}n en escenarios de alta densidad.} La
detecci\'{o}n en multitudes concentradas de individuos donde la oclusi\'{o}n mutua entre
personas es intensa y por ende requiere m\'{e}tricas y protocolos de
evaluaci\'{o}n m\'{a}s espec\'{i}ficos~\cite{E28, E09}.

\begin{hallazgo}[]
  La necesidad de datasets combinados y la demanda nesesitada de modelos unificados
  capaces de afrontar oclusi\'{o}n y baja visibilidad de forma
  simult\'{a}nea constituyen las brechas m\'{a}s frecuentemente
  citadas, apareciendo en el 53\,\% de los estudios primarios
  revisados~\cite{E31, E07, E17, E20, E26}.
\end{hallazgo}

\subsection{Mapa de Evidencia}

\begin{table}[H]
\centering
\caption{Mapa de evidencia: t\'{e}cnicas vs.\ preguntas de
investigaci\'{o}n.}
\label{tab:mapa}
\begin{tabular}{p{5.5cm}cccccc}
\toprule
\textbf{T\'{e}cnica / Enfoque} &
\textbf{RQ1} & \textbf{RQ2} & \textbf{RQ3} &
\textbf{RQ4} & \textbf{RQ5} & \textbf{RQ6} \\
\midrule
Fusi\'{o}n multiespectral RGB-IR/t\'{e}rmico
  & $\checkmark$ & $\checkmark$ & $\checkmark$ & $\checkmark$ & $\checkmark$ & $\checkmark$ \\
Variantes YOLO (YOLOv8, NP-YOLO, PYOLO)
  & $\checkmark$ & {--} & $\checkmark$ & $\checkmark$ & $\checkmark$ & $\checkmark$ \\
M\'{o}dulos de atenci\'{o}n espacial/canal
  & $\checkmark$ & $\checkmark$ & $\checkmark$ & {--} & $\checkmark$ & $\checkmark$ \\
Fusi\'{o}n multi-vista (c\'{a}maras m\'{u}ltiples)
  & $\checkmark$ & $\checkmark$ & $\checkmark$ & {--} & $\checkmark$ & $\checkmark$ \\
Transformers DETR (RC-DETR, LF-DETR)
  & $\checkmark$ & $\checkmark$ & $\checkmark$ & {--} & $\checkmark$ & $\checkmark$ \\
Estimaci\'{o}n de pose / esqueleto
  & $\checkmark$ & $\checkmark$ & {--} & {--} & $\checkmark$ & $\checkmark$ \\
NMS adaptado para multitudes
  & $\checkmark$ & $\checkmark$ & $\checkmark$ & $\checkmark$ & $\checkmark$ & {--} \\
Mejora de imagen previa al detector
  & $\checkmark$ & {--} & $\checkmark$ & {--} & $\checkmark$ & $\checkmark$ \\
\bottomrule
\end{tabular}
\end{table}

%=============================================================
\section{Discusi\'{o}n}
\label{sec:discusion}
%=============================================================

\subsection{Interpretaci\'{o}n de los Resultados}

Los hallazgos posteriores tras la realizaci\'{o}n de esta revisi\'{o}n son
consecuentes con la tendencia general documentada en trabajos previos
sobre detecci\'{o}n peatonal~\cite{Brunetti2018, Dollar2012}: la
calidad del detector usado depende cr\'{i}ticamente de la riqueza y
disponibilidad de informaci\'{o}n visual obetenida. Cuando la iluminaci\'{o}n o
la visibilidad se ven afectadas, la incorporaci\'{o}n de canales
espectrales añadidos aparece como la estrategia m\'{a}s efectiva
ya que proporciona una se\~{n}al complementaria no degradada por las
mismas condiciones que afectan al canal visible trabajado. La predominancia de
variantes de YOLO refleja la presi\'{o}n industrial por la nesesidad de soluciones
capaces de operar en tiempo real con hardware que sea  eficiente y a la ves accesible.
Arquitecturas como NP-YOLO~\cite{E12} y PYOLO~\cite{E18} demuestran
que es factible adaptar detectores de una etapa para condiciones
adversas sin cambiar la velocidad de inferencia, lo que los hace
atractivos para despliegues en c\'{a}maras IP con capacidad de
c\'{o}mputo embebido.

La distribuci\'{o}n de estudios por base de datos muestra que ACM
Digital Library y Scopus aportan conjuntamente el 57{,}1\,\% de los
estudios seleccionados, lo cual confirma que las conferencias
especializadas en visi\'{o}n por computadora y multimedia constituyen
una fuente altamente relevante de difusi\'{o}n de avances en este campo de la tecnologia.
ScienceDirect concentra el 34{,}3\,\% restante a trav\'{e}s de
revistas de Elsevier, mientras que Web of Science aporta estudios
complementarios de alto impacto.

\subsection{Brechas de Investigaci\'{o}n y Trabajo Futuro}

Las cinco brechas de conocimiento identificadas en la RQ6 configuran
una agenda de investigaci\'{o}n clara. La que posee mas urgencia es la
falta considerable de benchmarks que combinen simult\'{a}neamente
oclusi\'{o}n densa y condiciones meteorol\'{o}gicas dificiles. Los
conjuntos existentes (Caltech, KAIST, FLIR) abordan estas dimensiones
de forma separada, lo cual impide la validaci\'{o}n cruzada de modelos
que pretendan ser robustos ante ambas condiciones cuando estan en simultaneo. La segunda brecha,
vinculada a la adaptaci\'{o}n de dominio, cobra alta relevancia en
el contexto y situaciones del area de latinoamerica, donde las condiciones de infraestructura de hardware
tecnol\'{o}gico y los patrones de comportamiento peatonal pueden
diferir significativamente de los contextos estudiados en los
art\'{i}culos revisados conjuntamente.

\subsection{Amenazas a la Validez}

Las principales amenazas significativas a la validez de esta revisi\'{o}
se clasifican de acuerdo con~\cite{Wohlin2012}:

\begin{itemize}
  \item \textbf{Validez de la construcci\'{o}n:} la cadena de
        b\'{u}squeda tiene la posiblidad existente de no recuperar estudios que usen terminolog\'{i}a
        alternativa para describir oclusi\'{o}n o baja visibilidad en las camaras usadas. Esta
        amenaza se disminuy\'{o} mediante la inclusi\'{o}n de
        sin\'{o}nimos y la validaci\'{o}n iterativa de la cadena contra
        un conjunto de art\'{i}culos conocidos.
  \item \textbf{Validez interna:} la selecci\'{o}n y an\'{a}lisis de
        calidad de los estudios fueron realizadas en conjunto por el equipo
        investigador, lo que introduce riesgo de sesgo de
        inclusi\'{o}n. Se mitig\'{o} mediante la adecuada aplicaci\'{o}n
        sistem\'{a}tica de los criterios definidos en el protocolo.
  \item \textbf{Validez externa:} los resultados se limitan al
        per\'{i}odo 2022--2026 y a las cuatro bases de datos
        consultadas; estudios en otras fuentes o en idiomas distintos
        del ingl\'{e}s y espa\~{n}ol quedaron excluidos.
  \item \textbf{Validez de conclusi\'{o}n:} la heterogeneidad de
        m\'{e}tricas y grupos de datos reportados dificulta significativamente la
        comparaci\'{o}n quantitativa directa entre los modelos, por lo que
        la s\'{i}ntesis del tema se realiz\'{o} de manera narrativa.
\end{itemize}

%=============================================================
\section{Conclusiones}
%=============================================================

Esta Revisi\'{o}n Sistem\'{a}tica de Literatura sintetiz\'{o} la
evidencia publicada entre 2022 y 2026 sobre los modelos existentes para la
detecci\'{o}n autom\'{a}tica de personas bajo condiciones considerables
de oclusi\'{o}n y baja visibilidad en sistemas de videovigilancia de individuos, a
partir de 35 estudios primarios identificados en ACM Digital Library,
Scopus, ScienceDirect y Web of Science, se siguio el protocolo de
Kitchenham y Charters.

Respecto a los modelos existentes (RQ1), la fusi\'{o}n multiespectral
RGB-infrarrojo/\allowbreak t\'{e}rmico y las redes con m\'{o}dulos de atenci\'{o}n
constituyen los enfoques mas dominantes, seguidos por variantes compactas
de YOLO y estrategias  principalmente de multi-vista basadas en transformers. En cuanto a
las t\'{e}cnicas de manejo de la oclusi\'{o}n (RQ2), la fusi\'{o}n de
perspectivas m\'{u}ltiples a traves de varias camaras y las tecnicas de atenci\'{o}n cruzada
reportan las mejoras m\'{a}s consistentes. Las m\'{e}tricas m\'{a}s
utilizadas (RQ3) son mAP, tasa de error logar\'{i}tmica y FPS, mientras
que los benchmarks de mas uso (RQ4) son KAIST, FLIR Thermal, Caltech
y CityPersons. Las limitaciones predominantes (RQ5) incluyen la
dependencia constante de hardware especializado, la fragilidad ante oclusi\'{o}n
severa y la brecha de generalizaci\'{o}n. Las principales brechas (RQ6)
nos señalan la necesidad de datasets combinados y modelos unificados
capaces de gestionar ambas condiciones de forma simult\'{a}nea.

Las contribuciones originales de esta revisi\'{o}n son: (i)~una
taxonom\'{i}a general de los enfoques de detecci\'{o}n bajo condiciones de
visi\'{o}n adversas acotada al per\'{i}odo 2022--2026 con cobertura
de cuatro bases de datos indexadas; (ii)~un mapa de evidencia que
cruza enfoques con preguntas de investigaci\'{o}n; y (iii)~la
iden\-ti\-fi\-ca\-ci-\'{o}n sistem\'{a}tica de cinco vac\'{i}os de conocimiento
que orientan la agenda de investigaci\'{o}n futura en ambientes de
videovigilancia inteligente.

Para la pr\'{a}ctica, los resultados sugieren que los equipos de
desarrollo de sistemas de vigilancia deben darle prioridad a las arquitecturas de
fusi\'{o}n multiespectral cuando el presupuesto lo permita, y
m\'{o}dulos de atenci\'{o}n como alternativa de bajo coste ante
hardware restringido y limitado. Como l\'{i}neas de investigaci\'{o}n futura se
propone: (1)~desarrollar un benchmark estandarizado que combine
oclusi\'{o}n y condiciones meteorol\'{o}gicas adversamente dificiles de leer; (2)~explorar
t\'{e}cnicas de adaptaci\'{o}n de dominio para mejorar la
generalizaci\'{o}n; y (3)~dise\~{n}ar arquitecturas unificadas capaces
de gestionar simult\'{a}neamente ambas condiciones con una eficiencia
suficientemente alta para su despliegue y uso en dispositivos embebidos.

%=============================================================
\section*{Declaraciones}
%=============================================================

\noindent\textbf{Conflicto de intereses:} Los autores declaran no tener
conflicto de intereses relacionado con esta investigaci\'{o}n.

\noindent\textbf{Financiamiento:} Esta investigaci\'{o}n no cont\'{o}
con financiamiento externo. Fue desarrollada en el marco del curso de
Investigaci\'{o}n~III de la Escuela Profesional de Ingenier\'{i}a de
Sistemas de la Universidad Peruana Uni\'{o}n, Juliaca, Per\'{u}.

\noindent\textbf{Contribuciones de autor\'{i}a} (CRediT):\\
Celia Patricia Apaza Hilasaca: Conceptualizaci\'{o}n, proceso
metodol\'{o}gico, redacci\'{o}n del borrador original.\\
Brayan Raul Condori Quispe: Curaci\'{o}n de datos, b\'{u}squeda y
selecci\'{o}n de estudios, visualizaci\'{o}n.\\
Jhan Logan Ramos Quispe: Revisi\'{o}n y edici\'{o}n, evaluaci\'{o}n de
calidad, s\'{i}ntesis de resultados.

%=============================================================
\printbibliography[title={Referencias}, heading=bibintoc]
%=============================================================

%=============================================================
\appendix
%=============================================================

\section{Protocolo de Revisi\'{o}n Sistem\'{a}tica}

\begin{itemize}
  \item \textbf{ACM Digital Library:}
        [[Title: ``person detection''] OR [Title: ``pedestrian
        detection''] OR [Title: ``human detection'']] AND
        [[Abstract: ``surveillance''] OR [Abstract: ``occlusion''] OR
        [Abstract: ``low visibility''] OR [Abstract: ``infrared''] OR
        [Abstract: ``nighttime'']] AND Publication Date:
        (01/01/2022 TO 12/31/2026)
  \item \textbf{Scopus:} TITLE-ABS-KEY((``person detection'' OR
        ``pedestrian detection'' OR ``human detection'' OR ``people
        detection'') AND (``video surveillance'' OR ``surveillance'' OR
        ``CCTV'' OR ``security camera'' OR ``monitoring'') AND
        (``occlusion'' OR ``low visibility'' OR ``low illumination'' OR
        ``nighttime'' OR ``fog'' OR ``rain'') AND (``deep learning'' OR
        ``neural network'' OR ``convolutional'' OR ``object detection''
        OR ``feature extraction'')) AND PUBYEAR $>$ 2021 AND DOCTYPE
        IN (``ar'', ``cp'') AND LANGUAGE IN (``English'',
        `Spanish'')
  \item \textbf{ScienceDirect:} (``pedestrian detection'' OR ``person
        detection'') AND (``video surveillance'' OR ``CCTV'') AND
        (``occlusion'' OR ``low visibility'') AND (``deep learning'' OR
        ``neural network'')
  \item \textbf{Web of Science:} TS=((``person detection'' OR
        ``pedestrian detection'' OR ``human detection'') AND
        (``video surveillance'' OR ``CCTV'') AND (``occlusion'' OR
        ``low visibility'' OR ``low illumination'') AND (``deep
        learning'' OR ``neural network'') AND (``real-time'' OR
        ``accuracy'' OR ``performance'')) AND PY=(2022--2026)
\end{itemize}

\section{Lista Completa de Estudios Primarios}

\begin{singlespace}
\small
\begin{longtable}{>{\centering\arraybackslash}p{1cm} >{\raggedright\arraybackslash}p{2.5cm} >{\raggedright\arraybackslash}p{5.8cm} >{\centering\arraybackslash}p{1.2cm} >{\raggedright\arraybackslash}p{2.8cm}}
\caption{Estudios primarios incluidos en la revisi\'{o}n.
\label{tab:ap_estudios}}\\
\toprule
\textbf{ID} & \textbf{Autores} & \textbf{T\'{i}tulo} &
\textbf{A\~{n}o} & \textbf{Fuente} \\
\midrule
\endfirsthead
\midrule
\textbf{ID} & \textbf{Autores} & \textbf{T\'{i}tulo} &
\textbf{A\~{n}o} & \textbf{Fuente} \\
\midrule
\endhead\endfoot\bottomrule\endlastfoot
E01 & Niu et al.       & Physics-based adversarial attack on NIR human detector & 2023 & ACM MM \\
E02 & Li et al.        & SAM-guided semantic knowledge fusion VIS-IR & 2025 & ACM ICMR \\
E03 & Zhang et al.     & Occluded pedestrian detection cross-modal attention & 2024 & ACM ICIGP \\
E04 & Wang et al.      & Thermal-RGB fusion adverse weather surveillance & 2023 & ACM IVSP \\
E05 & Liu et al.       & Attention-based occlusion person detection crowds & 2022 & ACM MMAsia \\
E06 & Park et al.      & Real-time pedestrian low-light multispectral CCTV & 2024 & ACM IVSP \\
E07 & Chen et al.      & Occluded pedestrian visible-infrared feature align & 2024 & ACM ICMVA \\
E08 & Zhao et al.      & Nighttime person detection infrared-guided & 2025 & ACM IVSP \\
E09 & Li et al.        & Multi-view pedestrian detection residual mask fusion & 2026 & Future Gen. Comput. Syst. \\
E10 & Qi et al.        & LF-DETR aerial RGB-infrared pedestrian detection & 2026 & Remote Sensing \\
E11 & Tang et al.      & Adaptive pedestrian detection occlusion-aware & 2026 & Pattern Recognit. \\
E12 & Huang \& Liu     & NP-YOLO nighttime pedestrian video surveillance & 2025 & CVAA Conf. \\
E13 & Zhang et al.     & Enhanced SlowFast nighttime port surveillance & 2025 & Conf. Proc. \\
E14 & Li, W.           & Enhanced thermal imaging human detection & 2025 & IAECST Conf. \\
E15 & Connan et al.    & Real-time gun and person detection surveillance & 2025 & ACM SAC \\
E16 & Mudavath \& Niranjan & Person detection thermal images YOLOv8 & 2025 & Earth Sci. Inform. \\
E17 & He et al.        & Cross-modality differential multispectral detection & 2025 & Conf. Proc. \\
E18 & Hu et al.        & PYOLO plus YOLO occlusion pedestrian detection & 2025 & CCDC Conf. \\
E19 & Gupta \& Batra   & Multispectral pedestrian CNN channel fusion & 2025 & WCONF Conf. \\
E20 & Vu et al.        & Block-recurrent transformer thermal imaging & 2025 & Mach. Vis. Appl. \\
E21 & Dai et al.       & Multi-exposure YOLO nighttime pedestrian & 2026 & Signal Process.: Image Commun. \\
E22 & Guo et al.       & Pedestrian detection vision-language semantics & 2025 & Pattern Recognit. Lett. \\
E23 & Zhang et al.     & Enhanced lightweight YOLOv8s small-scale & 2025 & Digit. Signal Process. \\
E24 & Andrade et al.   & SUSAN violence detection surveillance & 2025 & Expert Syst. Appl. \\
E25 & Huang et al.     & Ms-VLPD multi-scale pedestrian detection & 2025 & Expert Syst. Appl. \\
E26 & Shen et al.      & EAFF-Net dual-modality pedestrian detection & 2025 & Infrared Phys. Technol. \\
E27 & Kim et al.       & Skeleton AI fall detection occlusions & 2025 & Autom. Constr. \\
E28 & Gao et al.       & Ranking-based DETRs crowded pedestrian & 2025 & Neurocomputing \\
E29 & Gao et al.       & RC-DETR crowded pedestrian contrastive & 2025 & Neural Networks \\
E30 & Qiu et al.       & PPM boolean optimizer multi-view pedestrian & 2024 & Pattern Recognit. \\
E31 & Deng et al.      & Lightweight cross-modal multispectral pedestrian & 2024 & Comput. Mater. Contin. \\
E32 & Tang et al.      & PFEL-Net lightweight multi-scale pedestrian & 2024 & J. King Saud Univ. \\
E33 & Kolluri et al.   & Surveillance video pedestrian transport ML & 2026 & Discov. Appl. Sci. \\
E34 & Akshatha et al.  & Human detection aerial thermal Faster R-CNN + SSD & 2022 & Electronics \\
E35 & Dao et al.       & Noise-aware DL multispectral pedestrian & 2022 & J. Electron. Imaging \\
\end{longtable}
\end{singlespace}

\section{Formularios de Extracci\'{o}n de Datos}

Los formularios de extracci\'{o}n individuales de cada estudio
primario recogen los campos definidos en la
Tabla~\ref{tab:extraccion}: identificador, t\'{i}tulo, autor(es),
a\~{n}o de publicaci\'{o}n, fuente, tipo de publicaci\'{o}n, base de
datos de origen, modelo o t\'{e}cnica de detecci\'{o}n utilizado,
tipo de oclusi\'{o}n abordada, condici\'{o}n de visibilidad,
m\'{e}tricas de evaluaci\'{o}n, conjunto de datos empleado,
resultados principales, limitaciones reportadas, vac\'{i}os y trabajo
futuro, dominio de aplicaci\'{o}n y puntuaci\'{o}n de calidad (CQI).
Los formularios completos est\'{a}n disponibles como material
suplementario.

\end{document}